% Generic one-page resume template (RenderCV-style layout).
%
% This is the structural starting point for every base resume. The /extract-profile
% skill copies this file to base/<role>/main.tex and fills it from your real resume.
% You can also copy it by hand:  cp base/_template/main.tex base/<role>/main.tex
%
% Replace every placeholder marked <LIKE THIS>. Keep the layout — it is tuned to fit
% one letter-size page. Build with:  latexmk -pdf -interaction=nonstopmode main.tex

\documentclass[10pt, letterpaper]{article}

% Packages:
\usepackage[
    ignoreheadfoot, % set margins without considering header and footer
    top=0.3 cm, % seperation between body and page edge from the top
    bottom=0.3 cm, % seperation between body and page edge from the bottom
    left=1.5 cm, % seperation between body and page edge from the left
    right=1.5 cm, % seperation between body and page edge from the right
    footskip=1.0 cm, % seperation between body and footer
    % showframe % for debugging
]{geometry} % for adjusting page geometry
\usepackage{titlesec} % for customizing section titles
\usepackage{tabularx} % for making tables with fixed width columns
\usepackage{array} % tabularx requires this
\usepackage[dvipsnames]{xcolor} % for coloring text
\definecolor{primaryColor}{RGB}{0, 0, 0} % define primary color
\usepackage{enumitem} % for customizing lists
\usepackage{fontawesome5} % for using icons
\usepackage{amsmath} % for math
\usepackage[
    pdftitle={Your Name's Resume},
    pdfauthor={Your Name},
    pdfcreator={LaTeX with RenderCV},
    colorlinks=true,
    urlcolor=primaryColor
]{hyperref} % for links, metadata and bookmarks
\usepackage[pscoord]{eso-pic} % for floating text on the page
\usepackage{calc} % for calculating lengths
\usepackage{bookmark} % for bookmarks
\usepackage{lastpage} % for getting the total number of pages
\usepackage{changepage} % for one column entries (adjustwidth environment)
\usepackage{paracol} % for two and three column entries
\usepackage{ifthen} % for conditional statements
\usepackage{needspace} % for avoiding page brake right after the section title
\usepackage{iftex} % check if engine is pdflatex, xetex or luatex

% Ensure that generate pdf is machine readable/ATS parsable:
\ifPDFTeX
    \input{glyphtounicode}
    \pdfgentounicode=1
    \usepackage[T1]{fontenc}
    \usepackage[utf8]{inputenc}
    \usepackage{lmodern}
\fi

\usepackage{charter}

% Some settings:
\raggedright
\AtBeginEnvironment{adjustwidth}{\partopsep0pt} % remove space before adjustwidth environment
\pagestyle{empty} % no header or footer
\setcounter{secnumdepth}{0} % no section numbering
\setlength{\parindent}{0pt} % no indentation
\setlength{\topskip}{0pt} % no top skip
\setlength{\columnsep}{0.15cm} % set column seperation
\pagenumbering{gobble} % no page numbering

\titleformat{\section}{\needspace{4\baselineskip}\bfseries\large}{}{0pt}{}[\vspace{1pt}\titlerule]

\titlespacing{\section}{
    % left space:
    -1pt
}{
    % top space:
    0.15 cm
}{
    % bottom space:
    0.15 cm
} % section title spacing

\renewcommand\labelitemi{$\vcenter{\hbox{\small$\bullet$}}$} % custom bullet points
\newenvironment{highlights}{
    \begin{itemize}[
        topsep=0.10 cm,
        parsep=0.10 cm,
        partopsep=0pt,
        itemsep=0pt,
        leftmargin=0 cm + 10pt
    ]
}{
    \end{itemize}
} % new environment for highlights


\newenvironment{highlightsforbulletentries}{
    \begin{itemize}[
        topsep=0.10 cm,
        parsep=0.10 cm,
        partopsep=0pt,
        itemsep=0pt,
        leftmargin=10pt
    ]
}{
    \end{itemize}
} % new environment for highlights for bullet entries

\newenvironment{onecolentry}{
    \begin{adjustwidth}{
        0 cm + 0.00001 cm
    }{
        0 cm + 0.00001 cm
    }
}{
    \end{adjustwidth}
} % new environment for one column entries

\newenvironment{twocolentry}[2][]{
    \onecolentry
    \def\secondColumn{#2}
    \setcolumnwidth{\fill, 4.5 cm}
    \begin{paracol}{2}
}{
    \switchcolumn \raggedleft \secondColumn
    \end{paracol}
    \endonecolentry
} % new environment for two column entries

\newenvironment{twocolumnentry}[2][]{
    \onecolentry
    \def\secondColumn{#2}
    \setcolumnwidth{\fill, 6 cm}
    \begin{paracol}{2}
}{
    \switchcolumn \raggedleft \secondColumn
    \end{paracol}
    \endonecolentry
}

\newenvironment{threecolentry}[3][]{
    \onecolentry
    \def\thirdColumn{#3}
    \setcolumnwidth{, \fill, 4.5 cm}
    \begin{paracol}{3}
    {\raggedright #2} \switchcolumn
}{
    \switchcolumn \raggedleft \thirdColumn
    \end{paracol}
    \endonecolentry
} % new environment for three column entries

\newenvironment{header}{
    \setlength{\topsep}{0pt}\par\kern\topsep\centering\linespread{1.5}
}{
    \par\kern\topsep
} % new environment for the header

% save the original href command in a new command:
\let\hrefWithoutArrow\href

\begin{document}
    \newcommand{\AND}{\unskip
        \cleaders\copy\ANDbox\hskip\wd\ANDbox
        \ignorespaces
    }
    \newsavebox\ANDbox
    \sbox\ANDbox{$|$}

    \begin{header}
        \fontsize{20 pt}{20 pt}\selectfont Your Name

        \normalsize
        \kern 5.0 pt%
        \mbox{Role Title}%
        \kern 5.0 pt%
        \AND%
        \kern 5.0 pt%
        \mbox{your.email@example.com}%
        \kern 5.0 pt%
        \AND%
        \kern 5.0 pt%
        \mbox{(000)-000-0000}%
        \kern 5.0 pt%
        \AND%
        \kern 5.0 pt%
        \mbox{\hrefWithoutArrow{https://github.com/yourhandle}{GitHub}}%
    \end{header}

    \section{Summary}
    Two- to three-sentence summary of who you are professionally, your strongest
    skills, and the kind of work you do. Replace this with content tailored from
    your real resume.

    \section{Professional Experience}
        \begin{twocolentry}{
            City, ST
        }
            \textbf{Most Recent Employer}\end{twocolentry}
        \vspace{0.05 cm}
        \begin{twocolentry}{
            Mon YYYY -- Mon YYYY
        }
            \textit{Job Title}\end{twocolentry}
        \vspace{0.05 cm}
        \begin{onecolentry}
            \begin{highlights}
                \item Accomplishment bullet: what you did, the tools you used, and the
                measurable result.
                \item Another accomplishment bullet, quantified where possible.
            \end{highlights}
        \end{onecolentry}
        \vspace{0.10 cm}

        \begin{twocolentry}{
            City, ST
        }
            \textbf{Previous Employer}\end{twocolentry}
        \vspace{0.05 cm}
        \begin{twocolentry}{
            Mon YYYY -- Mon YYYY
        }
            \textit{Job Title}\end{twocolentry}
        \vspace{0.05 cm}
        \begin{onecolentry}
            \begin{highlights}
                \item Accomplishment bullet.
                \item Accomplishment bullet.
            \end{highlights}
        \end{onecolentry}

    \section{Education}
        \begin{twocolentry}{
            Mon YYYY -- Mon YYYY
        }
            \textbf{University Name}, Degree and Field of Study\end{twocolentry}
        \begin{onecolentry}
            \begin{highlights}
                \begin{twocolentry}{
                    \textbf{GPA:} 0.0/4.0
                }
                    \textbf{Relevant coursework / focus:} list here
                \end{twocolentry}
            \end{highlights}
        \end{onecolentry}

    \section{Projects}
        \begin{twocolentry}{
        }
            \textbf{Project Name}\end{twocolentry}
        \vspace{0.10 cm}
        \begin{onecolentry}
            \begin{highlights}
                \item What the project did, the stack, and the outcome.
            \end{highlights}
        \end{onecolentry}

    \section{Technical Skills}
    \begin{onecolentry}
        \begin{highlightsforbulletentries}
            \item \textbf{Category:} skill, skill, skill
            \item \textbf{Category:} skill, skill, skill
            \item \textbf{Category:} skill, skill, skill
        \end{highlightsforbulletentries}
    \end{onecolentry}

\end{document}
